\section*{\Huge Conclusion}
This work, carried out in the framework of an internship during the final months of a Master's degree, has been dedicated to understanding and studying the evolution of quantum open systems under continuous measurement of their environment. We therefore described the fundamental theories of closed quantum systems up to the unconditional master equation describing the evolution of open quantum systems. In particular, the derivation of the Lindblad master equation presented in the first chapter was performed within the input-output formalism under the Born-Markov approximations. 

The heart of our study aimed on the impact of continuous monitoring of the environment of open quantum systems in the second chapter. This impact is specifically crystallized in the conditional trajectories depending on which process of measurement is made on the environment. We studied two of these processes in details, the first being described by a jump-like output signal (photo-detection), while the second one being characterized by a continuous diffusive signal (homodyne detection). The studied processes have been modeled with stochastic processes, namely a Poisson process and a Wiener process. We also extended this study to the theoretical expressions allowing to simulate trajectories, whose underlying stochastic differential equations (called stochastic master equations) are initially non-linear and challenging to solve. This has been made possible by deriving linear quantum trajectories and by deriving the Kraus operators formalism. Finally, we simulated some trajectories of a qubit and a harmonic oscillator under the continuous monitoring of their environment. These simulations clearly highlighted the validity of the theoretical results since the averaged trajectory is perfectly coincides with the unconditional trajectory (when the number of trajectories tends to infinity). This property explains precisely why the derivation of a specific stochastic master equation is called an ``unravelling'' of the Lindblad unconditional master equation.

Beyond the theoretical results established in this report, several directions can extend this work. While our study focused exclusively on a vacuum bath, a natural next step would be to incorporate thermal environments~\cite{W93}. Furthermore, introducing more realistic descriptions would bridge the gap toward experimental reality, such as accounting for imperfect measurement processes and detector inefficiencies~\cite{Albarelli}. In the same vein, the framework of homodyne detection can be generalized to the so-called ``heterodyne'' or even "general-dyne" detection schemes~\cite{Albarelli, Jacobs2006, Genoni14}. Beyond the standard systems considered here, the scientific community studied more complex system, e.g., a single nonlinear Kerr resonator whose Hamiltonian features quadratic driving terms~\cite{Bartolo2017}. Finally, the theory of continuous quantum measurement could be naturally integrated into the (huge) framework of quantum estimation and metrology~\cite{PARIS2009, Rossi2020}.